\chapter{Related Work}
\label{sec:related-work}

% 1-pager

Before we test out \openzl{}, we need to take a look at the current state of
compression of scientific data, and particularly its use (or lack thereof) in
high-performance computing (HPC). There has been a reasonable amount of research
into both lossless and lossy compression of scientific data. The structure of
this section will be as follows:
\begin{itemize}
\item Scientific compressors
\item HPC compression
\item Artificial intelligence (AI) compression
\item Lossy compressors
\item Hardware-accelerated compression
\end{itemize}

First, we will cover the current efforts towards compression of high-precision
floating point scientific data. Second, we will cover efforts of integrating
these compressors into HPC workloads. Third, briefly, we will cover compression
in AI-related environments. This is not \emph{directly} related to HPC, but

The last two sections (lossy compression and hardware acceleration) are not
directly relevant to my work. \openzl{} is lossless, so I will not
\emph{directly} compare it to lossy compressors. Lossy compression will,
generally, have better compression ratios than a lossless compressor at the cost
of some precision\cite{lu_understanding_2018}.


\section{Scientific Compression}
\label{sec:sci-compression}
Previous work has shown compression to be beneficial in HPC workloads.
\cite{ratanaworabhan_fast_2006} showed their compression methodology allowed for
a compression throughput rate that ``exceeds the throughput of Fast Ethernet
networks and of many hard disks.'' Additionally, the authors postulate ``the
time saved due to sending or storing the shorter compressed data stream is
likely to outweigh the runtime overhead introduced by performing the
compression'' \cite{ratanaworabhan_fast_2006}. However, that paper was 20 years
ago. In the interim, data sizes and network, storage and compute speeds have all
increased.


\subsection{Lossy Compression}
Lossy compressors such as SZ show compression ratios between
3.3:1--436:1\cite{di_fast_2016}. Even on the low end, this is better than
lossless compression schemes which generally only see ratios of less than
2.\cite{lu_understanding_2018}

So if lossy compression is so wonderful, why bother investigating OpenZL or
other lossless compressor? Lossy compression, by definition, loses data. In HPC
with high-precision floating-point data, precision loss may be unacceptable.
Precision loss may be acceptable as shown by \cite{di_fast_2016}
and~\cite{lu_understanding_2018}. However, this takes additional work and
verification to ensure the data loss has not affected the simulation output. In
this paper, I will focus on lossless compression because it can be implemented
with no loss to data accuracy.

\subsection{Hardware Acceleration}
The ``correct'' answer is to offload compression of data to specialized
hardware.\cite{li_performance_2026} Previous work
(\cite{li_accelerating_2024},~\cite{li_compression_2024}) shows NVIDIA BlueField
DPUs and its C-engine providing a significant speedup ($10\times$--$100\times$)
for compression time compared to its SoC. However, to make use of these DPUs
would require a not insignificant monetary cost; requiring purchasing DPUs for
each node in a cluster if one wanted to compress/decompress in transit. There is
the potential to buy fewer than $n$ DPUs for $n$ nodes, and compress on a DPU
and decompress on the CPU; that is outside the scope of this paper. To that end,
this paper will only focus on compression performed on the CPU\@. I will
investigate purely software solutions which come at no additional (monetary)
cost to implement.
